\section{Introdução}

A verificação da higiene de superfícies em linhas de produção é exigida em programas de segurança
de alimentos e de controle de contaminação. Os métodos de referência, baseados na coleta de
amostras de superfície \citep{iso18593} e no cultivo em placa, são confiáveis, mas o resultado
demora de um a vários dias, o que impede uma decisão no mesmo turno de produção.

Biossensores eletroquímicos de DNA (genossensores) detectam sequências específicas de
microrganismos, como trechos do gene do rRNA 16S, por meio de uma sonda de captura imobilizada no
eletrodo: a hibridização com o alvo altera a transferência de carga de uma sonda redox em solução
(por exemplo, o par ferri/ferrocianeto), o que aparece como mudança na corrente de pico do
voltamograma. Eletrodos impressos de carbono (SPCE) tornam a medição barata e portátil, mas o
sinal varia entre eletrodos, lotes e operadores, e a decisão por um limiar fixo de corrente tende a
falhar fora do laboratório.

Técnicas de aprendizado de máquina e quimiometria têm sido usadas para extrair mais informação de
sinais eletroquímicos, de PCA e PLS a florestas aleatórias, máquinas de vetores de suporte e redes
neurais \citep{puthongkham2021, bocan2025}. Duas dificuldades se repetem nesses trabalhos: os
conjuntos de dados são pequenos e heterogêneos, e a validação nem sempre impede que medições da
mesma amostra apareçam no treino e no teste, o que infla os resultados \citep{kapoor2023}. Métodos
que produzem descritores explícitos, como o SISSO \citep{ouyang2018}, são atraentes nesse cenário
porque funcionam com poucas amostras e dão uma fórmula interpretável.

Este trabalho descreve uma plataforma integrada, desenvolvida no projeto FIEP, com os seguintes
objetivos: (i) ler os voltamogramas diretamente dos arquivos do potenciostato e organizá-los em um
banco de dados rastreável, com controle de qualidade automático; (ii) comparar de forma sistemática
representações do sinal e modelos supervisionados, incluindo o SISSO, para classificar superfícies
e estimar a carga bacteriana; e (iii) disponibilizar o modelo escolhido em uma ferramenta de uso no
laboratório e nas empresas validadoras.
